\subsection{仿紧空间上的覆叠}

命题 8. --- 仿紧空间 (I, p. 69) 上的覆叠是仿紧空间。

先证明一个引理。

引理. --- 设 E 为分离的拓扑空间。设 E 有局部有限开覆盖 \((\mathrm{V}_{i})_{i\in\mathrm{I}}\)，使得对每个 \(i\in I\)，\(\overline{V_{i}}\) 都是仿紧空间。则空间 E 是仿紧的。

设 \((\mathrm{W}_{j})_{j\in\mathrm{J}}\) 为 E 的开覆盖；我们来证明存在 E 的局部有限开加细 \((\mathrm{A}_{k})_{k\in\mathrm{K}}\)，它比覆盖 \((\mathrm{W}_{j})_{j\in\mathrm{J}}\) 更细。对所有 \(i\in I\)，设 \((\mathrm{A}_{\ell}^{\prime})_{\ell\in\mathrm{K}_{i}}\) 为 \(\overline{V_{i}}\) 的由 \(\overline{V_{i}}\) 的开子集组成的局部有限覆盖，比覆盖 \((\mathrm{W}_{j}\cap\overline{\mathrm{V}_{i}})_{j\in\mathrm{J}}\) 更细。设 K 为族 \((\mathrm{K}_{i})_{i\in\mathrm{I}}\) 的和 (E, II, p. 30, def. 8)。对 K 的每个元素 \(k=(\ell,i)\)，令 \(A_{k}=A_{\ell}^{\prime}\cap V_{i}\)。于是 \(A_{k}\) 在 E 中是开的，且有 \(\bigcup_{k\in K}A_{k}=\bigcup_{i\in I}V_{i}=E\)；此外，对每个 \(k\in K\)，存在指标 \(j\in J\)，使得 \(A_{k}\subset W_{j}\)。于是族 \((\mathrm{A}_{k})_{k\in\mathrm{K}}\) 是 E 的比 \((\mathrm{W}_{j})_{j\in\mathrm{J}}\) 更细的开覆盖。剩下要证明族 \((\mathrm{A}_{k})_{k\in\mathrm{K}}\) 是局部有限的。设 \(x\in E\)；存在开邻域 U （x 的），它只与那些 \(V_{i}\) （其 i 属于 I 的有限子集 I'）相交。对每个 \(i\in I'\)，x 有开邻域 \(U_{i}\subset U\)，它只与有限多个开集 \(A_{k}\) （对 \(k\in K_{i}\times\{i\}\)）相交：若 x 不属于 \(\overline{V_{i}}\)，这是显然的；而若 x 属于 \(\overline{V_{i}}\)，这由覆盖 \((\mathrm{A}_{\ell}^{\prime})_{\ell\in\mathrm{K}_{i}}\) （\(\overline{V_{i}}\) 的）的局部有限性质得出。于是 \(V^{\prime}=\bigcap_{i\in I'}U_{i}\) 是 x 的开邻域，它只与诸 \(A_{k}\) （\(k\in K\)）中的有限多个相交，从而覆盖 \((\mathrm{A}_{k})_{k\in\mathrm{K}}\) 是局部有限的。

来证明该命题。设 E 为 B 的覆叠，用 p 表示其投影，并设空间 B 是仿紧的。设 \((\mathrm{A}_{i})_{i\in\mathrm{I}}\) 为 B 的局部有限开覆盖，使得对每个 \(i\in I\)，覆叠 E 在 \(A_{i}\) 上可平凡化。对所有 \(i\in I\)，设 \(F_{i}\) 为离散拓扑空间，并设 \(g_{i}\colon\overline{p}^{1}(\mathrm{A}_{i})\to\mathrm{A}_{i}\times\mathrm{F}_{i}\) 为 E 在 \(A_{i}\) 上的一个平凡化。设 \((\mathrm{B}_{i})_{i\in\mathrm{I}}\) 为 B 的开覆盖，使得对每个 \(i\in I\)，有 \(\overline{B_{i}}\subset A_{i}\) (TG, IX, p. 49, prop. 4 与 p. 48, cor. 1)。对每个 \(i\in I\)，令 \(V_{i}=\overline{p}^{1}(\mathrm{B}_{i})\)；我们有 \(\overline{V_{i}}\subset\overline{p}^{1}(\overline{\mathrm{B}_{i}})\subset\overline{p}^{1}(\mathrm{A}_{i})\) 以及 \(V_{i}=\overline{g_{i}}^{1}(\mathrm{B}_{i}\times\mathrm{F}_{i})\)，从而 \(\overline{V_{i}}\subset\overline{g_{i}}^{1}(\overline{\mathrm{B}_{i}}\times\mathrm{F}_{i})\)。由于 B 是仿紧的，\(\overline{B_{i}}\) 是仿紧的 (TG, I, p. 69, prop. 16)，且 \(\overline{B_{i}}\times F_{i}\) 是仿紧的 (TG, I, p. 70, prop. 18)。于是 \(\overline{g_{i}}^{1}(\overline{\mathrm{B}_{i}}\times\mathrm{F}_{i})\) 是仿紧的，从而 \(\overline{V_{i}}\) 也是仿紧的 (TG, I, p. 69, prop. 16)。族 \((\mathrm{V}_{i})_{i\in\mathrm{I}}\) 按构造是 E 的局部有限开覆盖。最后，空间 E 是分离的 (I, p. 26, remark 3)。因此它满足引理的假设，由此得出命题。

注. --- 可以证明，可度量化空间上的覆叠是可度量化的 (I, p. 145, exerc. 1)，而局部紧且无穷可数空间上的连通覆叠本身是局部紧且无穷可数的 (I, p. 145, exerc. 2)。

